\documentclass[12pt]{article}
\usepackage[utf8]{inputenc}
\usepackage[T1]{fontenc}
\usepackage{hyperref}
\title{Zapotec, San Agust\'in Mixtepec Language Resources for K-12 Curriculum Development}
\author{Compiled by Claude Sonnet 5.5 (claude-sonnet-5-5) for Andr\'as Kornai}
\date{}
\begin{document}
\maketitle

\section{Introduction}
San Agust\'in Mixtepec Zapotec (ISO 639-3 ztm, Glottocode sana1286) is an Otomanguean language of Oaxaca, Mexico, in the Sierra Sur group. It is severely endangered. Sources list about 59 speakers (a 1994 figure). The searches found almost no educational or documentary material.

\section{Primary Educational Resources}
None were found.

\section{Linguistic and Reference Materials}
\begin{itemize}
\item Wikipedia overview: \url{https://en.wikipedia.org/wiki/San_Agust%C3%ADn_Mixtepec_Zapotec}
\item Verbix endangered-languages page (describes the language as unwritten): \url{https://docs.verbix.com/EndangeredLanguages/SanAgustinMixtepecZapotec}
\item Endangered Languages Project entry: \url{https://endangeredlanguages.com/lang/3116}
\end{itemize}
No grammar or dictionary was found.

\section{Religious and Community Resources}
None were found.

\section{Digital and Audio Resources}
\begin{itemize}
\item Global Recordings Network page: \url{https://globalrecordings.net/en/language/ztm}. It states that no recordings are currently available.
\end{itemize}

\section{Curriculum Development Context}
With very few speakers and no known materials, the first need is urgent documentation: recordings, a word list and an orthography, done with the remaining speakers. School curriculum work would follow from that.
\end{document}